\subsection{Trace identities for block intersections}\label{subsec:block_intersection_trace}

\begin{lemma}[BNT block separation gives a finite blockwise product span]
    \label{lem:bnt_direct_sum_block_separation_word_span}
    \lean{MPSTensor.hasPairBlockSeparatingWords_threeBlock_of_blocksNotGaugePhaseEquiv_c1}
    \lean{MPSTensor.propBlockInjective_of_blocksNotGaugePhaseEquiv_directSum_c1}
    \lean{MPSTensor.wordTupleSpanTop_of_blocksNotGaugePhaseEquiv_directSum_selectors_c1}
    \lean{MPSTensor.exists_wordTupleSpanTop_of_blocksNotGaugePhaseEquiv_directSum_selectors_c1}
    \leanok
    \uses{lem:direct_sum_two_block_dimension_step,
        lem:pair_trace_separation_to_pairwise_block_separation,
        lem:block_separating_equations_from_pairwise_separation,
        lem:product_word_span_from_bnt_block_separation}
    Let \(A^1,\ldots,A^r\) be separated BNT blocks satisfying the
    irreducibility, left-canonicality, and normalized self-overlap hypotheses
    used in the direct-sum comparison. Fix \(L_0>0\). Assume that, for every
    block \(j\), the map
    \[
        \Gamma_{L_0}^{A^j}:M_{D_j}(\mathbb C)\to
        (\mathbb C^d)^{\otimes L_0},
        \qquad
        X\longmapsto
        \sum_{|u|=L_0}\tr(XA^j_u)\ket{u}
    \]
    is injective, and that the homogeneous word spans at the two source
    lengths are full:
    \[
        \operatorname{span}\{A^j_u:|u|=L_0+1\}=M_{D_j}(\mathbb C),
        \qquad
        \operatorname{span}\{A^j_v:|v|=3(L_0+1)\}=M_{D_j}(\mathbb C).
    \]
    Then for every ordered pair \(k\ne j\) there are coefficients
    \(c^{k,j}_\omega\) such that
    \[
        \sum_{|\omega|=3(L_0+1)} c^{k,j}_\omega A^k_\omega=I,
        \qquad
        \sum_{|\omega|=3(L_0+1)} c^{k,j}_\omega A^j_\omega=0.
    \]
    Consequently the simultaneous block-word tuples at length
    \[
        (L_0+1)+(r-1)3(L_0+1)
    \]
    span the full product algebra
    \[
        \prod_{j=1}^r M_{D_j}(\mathbb C).
    \]
\end{lemma}

\begin{proof}
    \leanok
    Lemma~\ref{lem:direct_sum_two_block_dimension_step} gives homogeneous
    trace separation at length \(3(L_0+1)\) for every ordered pair of distinct
    BNT blocks. Lemma~\ref{lem:pair_trace_separation_to_pairwise_block_separation}
    converts these trace-separation equations into the displayed pairwise
    block-separating equations. Multiplying the equations for the \(r-1\)
    blocks different from \(k\) gives a block-separating equation for \(k\), and
    Lemma~\ref{lem:product_word_span_from_bnt_block_separation} then supplies the
    displayed product span.
\end{proof}

\begin{lemma}[Blockwise matrix equality from trace pairings]
    \label{lem:block_matrices_equal_from_trace_pairings}
    \lean{MPSTensor.block_matrices_eq_of_wordTupleSpanTop_trace}
    \leanok
    \uses{def:blockwise_product_word_span}
    Let \(A^1,\ldots,A^g\) be block tensors whose length-\(m\) simultaneous word
    tuples span the product algebra. If, for every word \(w\) of length \(m\),
    two blockwise matrix families \(X_j\) and \(Y_j\) satisfy
    \[
        \sum_j\operatorname{tr}(X_jA^j_w)
        =
        \sum_j\operatorname{tr}(Y_jA^j_w)
    \]
    then \(X_j=Y_j\) for every block \(j\).
\end{lemma}

\begin{proof}
    \leanok
    \uses{def:blockwise_product_word_span}
    For every word \(w\) of length \(m\),
    \[
        \sum_j\operatorname{tr}\!\left((X_j-Y_j)A^j_w\right)
        =
        \sum_j\operatorname{tr}(X_jA^j_w)
        -
        \sum_j\operatorname{tr}(Y_jA^j_w)
        =
        0.
    \]
    Nondegeneracy of the product trace pairing, together with the spanning
    hypothesis, gives \(X_j-Y_j=0\) for every block \(j\).
\end{proof}

\begin{lemma}[Boundary identities from block-sum membership]
    \label{lem:blockwise_boundary_identities_from_left_trace_membership}
    \lean{MPSTensor.pgvwc07_boundary_identities_of_leftBoundaryComponent_mem_iSup}
    \leanok
    \uses{lem:block_matrices_equal_from_trace_pairings,
        def:blockwise_product_word_span,
        def:ground_space_parent}
    Let \(A^1,\ldots,A^g\) be block tensors whose length-\(m\) simultaneous word
    tuples span the product algebra. Let \(C^j_a\in M_{D_j}(\mathbb C)\) be
    matrices indexed by blocks \(j\) and physical letters \(a\). Suppose that,
    for every pair of boundary letters \(a,b\) and every word \(w\) of length
    \(m\), a vector \(\psi\) has the left-boundary trace form
    \[
        \psi(a,w,b)=
        \sum_j\operatorname{tr}\!\left(A^j_b C^j_a A^j_w\right)
    \]
    and if
    \[
        \psi\in\bigvee_j G_{m+2}(A^j),
    \]
    then there are matrices \(E_j\) such that, for every block \(j\) and all
    \(a,b\),
    \[
        A^j_bC^j_a=A^j_bE_jA^j_a
    \]
\end{lemma}

\begin{proof}
    \leanok
    \uses{lem:block_matrices_equal_from_trace_pairings,
        def:ground_space_parent,
        def:ground_space_map}
    Write \(\psi=\sum_j\phi_j\) with \(\phi_j\in G_{m+2}(A^j)\). For each
    \(j\), choose \(E_j\) whose image under the length-\(m+2\) ground-space map
    for \(A^j\) is \(\phi_j\). Comparing the coefficient of the word \(awb\)
    in the two expressions for \(\psi\) gives
    \[
        \sum_j\operatorname{tr}\!\left(A^j_b C^j_a A^j_w\right)
        =
        \sum_j\operatorname{tr}\!\left(A^j_bE_jA^j_aA^j_w\right).
    \]
    Lemma~\ref{lem:block_matrices_equal_from_trace_pairings} gives the displayed
    blockwise identities.
\end{proof}

\begin{lemma}[Blockwise boundary compatibility from traces]
    \label{lem:blockwise_boundary_compatibility_from_traces}
    \lean{MPSTensor.pgvwc07_blockwise_compatibility_of_trace_decomposition}
    \leanok
    \uses{lem:block_matrices_equal_from_trace_pairings}
    Let \(A^1,\ldots,A^g\) be block tensors with the following common
    word-span property: for a fixed \(m\), the simultaneous tuples
    \[
        w\longmapsto
        \left(A^1_w,\ldots,A^g_w\right),
        \qquad |w|=m,
    \]
    span the full product algebra
    \(\prod_j M_{D_j}(\mathbb C)\).
    Suppose that matrices \(C^j_a,D^j_b\in M_{D_j}(\mathbb C)\) satisfy,
    for all physical indices \(a,b\) and all words \(w\) of length \(m\),
    \[
        \sum_j \operatorname{tr}\!\left(A^j_b C^j_a A^j_w\right)
        =
        \sum_j \operatorname{tr}\!\left(D^j_b A^j_a A^j_w\right).
    \]
    Then
    \[
        A^j_bC^j_a=D^j_bA^j_a
    \]
    for every block \(j\) and all \(a,b\).
\end{lemma}

\begin{proof}
    \leanok
    \uses{lem:block_matrices_equal_from_trace_pairings}
    For fixed \(a,b\), set
    \[
        \Delta_j=A^j_bC^j_a-D^j_bA^j_a .
    \]
    The assumed trace equality says that
    \[
        \sum_j \operatorname{tr}\!\left(\Delta_jA^j_w\right)=0
    \]
    for every word \(w\) of length \(m\).  Since the tuples
    \((A^1_w,\ldots,A^g_w)\) span the product algebra, the product trace
    pairing is nondegenerate and hence every \(\Delta_j\) is zero.
\end{proof}

\begin{lemma}[Word-valued boundary compatibility from traces]
    \label{lem:blockwise_word_boundary_compatibility_from_traces}
    \lean{MPSTensor.pgvwc07_blockwise_word_compatibility_of_trace_decomposition,
        MPSTensor.pgvwc07_complementary_word_compatibility_of_trace_decomposition}
    \leanok
    \uses{lem:block_matrices_equal_from_trace_pairings}
    Let \(A^1,\ldots,A^g\) be block tensors with the common word-span property
    at length \(m\). Suppose that \(C^j_\rho \in M_{D_j}(\mathbb C)\) is
    indexed by a block \(j\) and a word \(\rho\) of some fixed length \(M\),
    and that \(D^j_\beta \in M_{D_j}(\mathbb C)\) is indexed by a block
    \(j\) and a word \(\beta\) of some fixed length \(K\). If, for every
    middle word \(w\) of length \(m\),
    \[
        \sum_j\operatorname{tr}\!\left(A^j_\beta C^j_\rho A^j_w\right)
        =
        \sum_j\operatorname{tr}\!\left(D^j_\beta A^j_\rho A^j_w\right),
    \]
    then
    \[
        A^j_\beta C^j_\rho=D^j_\beta A^j_\rho
    \]
    for every block \(j\) and all words \(\beta,\rho\).  In particular, taking
    \(D^j_\beta=X_jA^j_\beta\) gives
    \[
        A^j_\beta C^j_\rho=(X_jA^j_\beta)A^j_\rho .
    \]
\end{lemma}

\begin{proof}
    \leanok
    \uses{lem:block_matrices_equal_from_trace_pairings}
    Fix \(\beta\) and \(\rho\), and set
    \[
        \Delta_j=A^j_\beta C^j_\rho-D^j_\beta A^j_\rho .
    \]
    The trace equality says that
    \[
        \sum_j\operatorname{tr}(\Delta_jA^j_w)=0
    \]
    for every word \(w\) of length \(m\).  The same nondegenerate product
    trace-pairing argument as in
    Lemma~\ref{lem:blockwise_boundary_compatibility_from_traces} gives
    \(\Delta_j=0\) for every \(j\).  The specialization
    \(D^j_\beta=X_jA^j_\beta\) is immediate.
\end{proof}

\begin{lemma}[Normalized boundary matrix identities]
    \label{lem:normalized_boundary_matrix_identities}
    \lean{MPSTensor.pgvwc07_boundary_matrix_identities_of_indexed_compatibility}
    \lean{MPSTensor.pgvwc07_boundary_matrix_identities_of_compatibility}
    \lean{MPSTensor.pgvwc07_boundary_word_matrix_identities_of_compatibility}
    \lean{MPSTensor.pgvwc07_boundary_matrix_identities_of_two_index_compatibility}
    \lean{MPSTensor.pgvwc07_boundary_word_matrix_identities_of_two_length_compatibility}
    \leanok
    Let $(A_a)_{a\in I}$, $(C_a)_{a\in I}$, and $(D_b)_{b\in I}$ be matrices in
    $M_D(\mathbb C)$, indexed by a finite alphabet. If
    \[
        \sum_a A_a A_a^\dagger=\Id,
        \qquad
        A_b C_a=D_b A_a
    \]
    for all $a,b\in I$, and if
    \[
        E=\sum_a C_a A_a^\dagger,
    \]
    then
    \[
        D_b=A_bE,
        \qquad
        A_bC_a=A_bEA_a
    \]
    for all $a,b\in I$.

    In particular, the same calculation applies when the alphabet is the finite
    set of words of a fixed length and $A_\beta$ denotes
    $A_{\beta_1}\cdots A_{\beta_K}$.

    The same argument also applies when the normalized family is indexed by a
    finite set $J$, while $(A_b)_{b\in I}$ and $(D_b)_{b\in I}$ are indexed by
    a set $I$. If $(B_a)_{a\in J}$ and $(C_a)_{a\in J}$ are indexed by $J$,
    and if for all $a\in J$ and $b\in I$,
    \[
        \sum_a B_aB_a^\dagger=\Id,
        \qquad
        A_bC_a=D_bB_a
    \]
    then with
    \[
        E=\sum_a C_aB_a^\dagger
    \]
    one has
    \[
        D_b=A_bE,
        \qquad
        A_bC_a=A_bEB_a.
    \]
    This includes the case where $I$ and $J$ are word sets of different
    lengths.
\end{lemma}

\begin{proof}
    \leanok
    Multiply $D_b$ by the normalization:
    \[
        D_b
        =
        D_b\sum_a A_aA_a^\dagger
        =
        \sum_a D_bA_aA_a^\dagger
        =
        \sum_a A_bC_aA_a^\dagger
        =
        A_bE .
    \]
    Substituting this in $A_bC_a=D_bA_a$ gives
    $A_bC_a=A_bEA_a$.
    For two alphabets, the corresponding step is
    \[
        D_b
        =
        D_b\sum_a B_aB_a^\dagger
        =
        \sum_a D_bB_aB_a^\dagger
        =
        \sum_a A_bC_aB_a^\dagger
        =
        A_bE ,
    \]
    where $E=\sum_a C_aB_a^\dagger$.  Substituting this in
    $A_bC_a=D_bB_a$ gives $A_bC_a=A_bEB_a$.
\end{proof}

\begin{lemma}[Complementary-word boundary identities from the PGVWC comparison]
    \label{lem:complementary_word_boundary_identities_from_pgvwc}
    \lean{MPSTensor.pgvwc07_complementary_word_boundary_identities_of_compatibility}
    \leanok
    \uses{lem:normalized_boundary_matrix_identities}
    Let $A_\beta$ be the products over wrapped words of length $K$, let
    $A_\rho$ be the products over complementary words of length $M$, and let
    $X\in M_D(\mathbb C)$ be a matrix.  Suppose that the complementary-word
    products satisfy
    \[
        \sum_\rho A_\rho A_\rho^\dagger=\Id
    \]
    and that there are matrices $C_\rho$, indexed by complementary words, such
    that for every wrapped word $\beta$ and complementary word $\rho$,
    \[
        A_\beta C_\rho=(X A_\beta)A_\rho.
    \]
    Then, for every complementary word $\rho$, there is a matrix $E_\rho$ such
    that for every wrapped word $\beta$,
    \[
        (X A_\beta)A_\rho=A_\beta E_\rho.
    \]
\end{lemma}

\begin{proof}
    \leanok
    \uses{lem:normalized_boundary_matrix_identities}
    Apply Lemma~\ref{lem:normalized_boundary_matrix_identities} with
    $D_\beta=X A_\beta$. It gives
    \[
        A_\beta C_\rho=A_\beta E A_\rho
    \]
    for $E=\sum_\gamma C_\gamma A_\gamma^\dagger$. With $E_\rho=EA_\rho$,
    the compatibility hypothesis gives
    \[
        (X A_\beta)A_\rho
        = A_\beta C_\rho
        = A_\beta E A_\rho
        = A_\beta E_\rho.
    \]
\end{proof}

\begin{lemma}[A left-boundary summand is a ground-space vector]
    \label{lem:left_boundary_component_mem_ground_space}
    \lean{MPSTensor.pgvwc07LeftBoundaryComponent_mem_groundSpace}
    \leanok
    Let \(A\) be a block tensor, and let \(C_a,E\in M_D(\mathbb C)\)
    satisfy
    \[
        A_b C_a=A_bEA_a
    \]
    for all physical indices \(a,b\).  For a word
    \(\sigma=(\sigma_1,\ldots,\sigma_{n+2})\), define
    \[
        \alpha(\sigma)
        =
        \operatorname{tr}\!\left(
        A_{\sigma_{n+2}} C_{\sigma_1}
        A_{\sigma_2}\cdots A_{\sigma_{n+1}}
        \right).
    \]
    Then \(\alpha\in G_{n+2}(A)\).
\end{lemma}

\begin{proof}
    \leanok
    The boundary identity gives
    \[
        A_{\sigma_{n+2}} C_{\sigma_1}
        =
        A_{\sigma_{n+2}}EA_{\sigma_1}.
    \]
    Hence, by cyclicity of the trace,
    \[
        \alpha(\sigma)
        =
        \operatorname{tr}\!\left(
        A_{\sigma_1}A_{\sigma_2}\cdots A_{\sigma_{n+2}}E
        \right),
    \]
    which is the ground-space parametrization of \(E\).
\end{proof}

\begin{lemma}[Left-boundary summands are ground-space vectors]
    \label{lem:left_boundary_summands_mem_block_ground_spaces}
    \lean{MPSTensor.pgvwc07_sum_leftBoundaryComponents_mem_iSup_groundSpace}
    \leanok
    \uses{lem:left_boundary_component_mem_ground_space}
    Let \(A^1,\ldots,A^g\) be block tensors, and let
    \(C^j_a,E^j\in M_{D_j}(\mathbb C)\) satisfy
    \[
        A^j_b C^j_a=A^j_bE^jA^j_a
    \]
    for every block \(j\) and all physical indices \(a,b\).
    For a word \(\sigma=(\sigma_1,\ldots,\sigma_{n+2})\), define
    \[
        \alpha_j(\sigma)
        =
        \operatorname{tr}\!\left(
        A^j_{\sigma_{n+2}} C^j_{\sigma_1}
        A^j_{\sigma_2}\cdots A^j_{\sigma_{n+1}}
        \right).
    \]
    Then
    \[
        \sum_j \alpha_j
        \in
        \bigvee_j G_{n+2}(A^j).
    \]
\end{lemma}

\begin{proof}
    \leanok
    \uses{lem:left_boundary_component_mem_ground_space}
    Fix \(j\).  By the assumed boundary identity and cyclicity of the
    trace,
    \[
        \operatorname{tr}\!\left(
        A^j_{\sigma_{n+2}} C^j_{\sigma_1}
        A^j_{\sigma_2}\cdots A^j_{\sigma_{n+1}}
        \right)
        =
        \operatorname{tr}\!\left(
        A^j_{\sigma_1}A^j_{\sigma_2}\cdots
        A^j_{\sigma_{n+2}}E^j
        \right).
    \]
    Thus \(\alpha_j\) is the image of \(E^j\) under the parametrization
    of \(G_{n+2}(A^j)\).  Hence each \(\alpha_j\) lies in
    \(G_{n+2}(A^j)\), and the finite sum lies in the supremum of these
    subspaces.
\end{proof}

\begin{lemma}[Left-boundary trace decompositions lie in the block ground spaces]
    \label{lem:left_boundary_trace_decomposition_mem_block_ground_spaces}
    \lean{MPSTensor.pgvwc07_mem_iSup_groundSpace_of_trace_decomposition}
    \leanok
    Let \(A^1,\ldots,A^g\) be block tensors with common word-span separation at
    length \(n\), and suppose that
    \[
        \sum_a A^j_a A^{j\,\dagger}_a=\Id
    \]
    for every block \(j\).  Suppose a vector \(\psi\in(\mathbb C^d)^{\otimes(n+2)}\)
    has the left-boundary trace decomposition
    \[
        \psi=\sum_j\alpha_j,
        \qquad
        \alpha_j(i_1,\ldots,i_{n+2})
        =
        \operatorname{tr}\!\left(
        A^j_{i_{n+2}} C^j_{i_1}
        A^j_{i_2}\cdots A^j_{i_{n+1}}
        \right),
    \]
    and suppose that the two trace decompositions agree for every middle word:
    \[
        \sum_j
        \operatorname{tr}\!\left(
        A^j_b C^j_a A^j_w
        \right)
        =
        \sum_j
        \operatorname{tr}\!\left(
        D^j_b A^j_a A^j_w
        \right).
    \]
    Then
    \[
        \psi\in\bigvee_j G_{n+2}(A^j).
    \]
\end{lemma}

\begin{proof}
    \leanok
    The trace-decomposition equality and the common word-span hypothesis give
    \[
        A^j_bC^j_a=D^j_bA^j_a
    \]
    for every \(j,a,b\).  By the normalization,
    \[
        D^j_b
        =
        D^j_b\sum_a A^j_aA^{j\,\dagger}_a
        =
        \sum_a A^j_bC^j_aA^{j\,\dagger}_a
        =
        A^j_bE^j,
        \qquad
        E^j:=\sum_a C^j_aA^{j\,\dagger}_a.
    \]
    Hence \(A^j_bC^j_a=A^j_bE^jA^j_a\).  Lemma
    \ref{lem:left_boundary_summands_mem_block_ground_spaces} then gives
    \[
        \psi=\sum_j\alpha_j\in\bigvee_j G_{n+2}(A^j).
    \]
\end{proof}

\begin{lemma}[One-step block intersection from block restrictions]
    \label{lem:block_restrictions_mem_block_ground_spaces}
    \lean{MPSTensor.pgvwc07_mem_iSup_groundSpace_of_iSup_restrictions}
    \leanok
    \uses{lem:left_boundary_trace_decomposition_mem_block_ground_spaces}
    Let \(A^1,\ldots,A^g\) be block tensors with common word-span separation at
    length \(n\), and suppose
    \[
        \sum_a A^j_aA^{j\,\dagger}_a=\Id
    \]
    for every block \(j\).  If \(\psi\in(\mathbb C^d)^{\otimes(n+2)}\) satisfies
    \[
        \psi(a,-)\in\bigvee_jG_{n+1}(A^j),
        \qquad
        \psi(-,b)\in\bigvee_jG_{n+1}(A^j)
    \]
    for every physical index \(a,b\), then
    \[
        \psi\in\bigvee_jG_{n+2}(A^j).
    \]
\end{lemma}

\begin{proof}
    \leanok
    Choose finite decompositions of the fixed-boundary vectors:
    \[
        \psi(a,i_2,\ldots,i_{n+2})
        =
        \sum_j
        \operatorname{tr}\!\left(
        A^j_{i_2}\cdots A^j_{i_{n+2}}C^j_a
        \right),
    \]
    and
    \[
        \psi(i_1,\ldots,i_{n+1},b)
        =
        \sum_j
        \operatorname{tr}\!\left(
        A^j_{i_1}\cdots A^j_{i_{n+1}}D^j_b
        \right).
    \]
    Evaluating these two decompositions on
    \((a,w,b)\), where \(w\) has length \(n\), gives
    \[
        \sum_j\operatorname{tr}\!\left(A^j_bC^j_aA^j_w\right)
        =
        \sum_j\operatorname{tr}\!\left(D^j_bA^j_aA^j_w\right).
    \]
    The first decomposition also gives
    \[
        \psi=\sum_j\alpha_j,\qquad
        \alpha_j(i_1,\ldots,i_{n+2})
        =
        \operatorname{tr}\!\left(
        A^j_{i_{n+2}}C^j_{i_1}
        A^j_{i_2}\cdots A^j_{i_{n+1}}
        \right).
    \]
    Lemma~\ref{lem:left_boundary_trace_decomposition_mem_block_ground_spaces}
    therefore gives
    \[
        \psi\in\bigvee_jG_{n+2}(A^j).
    \]
\end{proof}

\begin{lemma}[One-step block intersection characterization]
    \label{lem:block_restrictions_iff_block_ground_spaces}
    \lean{MPSTensor.pgvwc07_mem_iSup_groundSpace_iff_iSup_restrictions}
    \leanok
    \uses{lem:block_restrictions_mem_block_ground_spaces}
    Under the hypotheses of
    Lemma~\ref{lem:block_restrictions_mem_block_ground_spaces}, an
    \((n+2)\)-site vector satisfies
    \[
        \psi\in\bigvee_jG_{n+2}(A^j)
    \]
    if and only if
    \[
        \psi(a,-)\in\bigvee_jG_{n+1}(A^j),
        \qquad
        \psi(-,b)\in\bigvee_jG_{n+1}(A^j)
    \]
    for every physical index \(a,b\).
\end{lemma}

\begin{proof}
    \leanok
    The reverse implication is
    Lemma~\ref{lem:block_restrictions_mem_block_ground_spaces}.  For the
    forward implication, write
    \[
        \psi=\sum_j\gamma_j,\qquad \gamma_j\in G_{n+2}(A^j).
    \]
    Each \(\gamma_j\) has fixed-boundary restrictions in \(G_{n+1}(A^j)\):
    \[
        \gamma_j(a,-)\in G_{n+1}(A^j),
        \qquad
        \gamma_j(-,b)\in G_{n+1}(A^j).
    \]
    Taking the finite sum over \(j\) gives the two restrictions of \(\psi\) in
    \(\bigvee_jG_{n+1}(A^j)\).
\end{proof}

\begin{lemma}[Subspace form of the one-step block intersection]
    \label{lem:block_restriction_intersection_subspace}
    \lean{MPSTensor.pgvwc07_iSup_groundSpace_eq_restriction_intersection}
    \leanok
    \uses{lem:block_restrictions_iff_block_ground_spaces}
    Under the hypotheses of
    Lemma~\ref{lem:block_restrictions_mem_block_ground_spaces}, set
    \[
        S_n:=\bigvee_jG_{n+1}(A^j).
    \]
    Then
    \[
        \left\{\psi:\forall b,\ \psi(-,b)\in S_n\right\}
        \cap
        \left\{\psi:\forall a,\ \psi(a,-)\in S_n\right\}
        =
        \bigvee_jG_{n+2}(A^j).
    \]
\end{lemma}

\begin{proof}
    \leanok
    Lemma~\ref{lem:block_restrictions_iff_block_ground_spaces} gives the
    displayed equivalence for each vector \(\psi\).  Interpreting the two sides
    of that equivalence as subspaces gives the stated equality.
\end{proof}
